\section{Introducción}


Antes de ahondar en la propuesta de solución a la problemática que guía este proyecto, resulta importante tener en consideración la forma en la que esta ha sido atacada previamente y cómo el reciente desarrollo tecnológico ha permitido que tareas como el reconocimiento facial puedan implementarse en distintos ámbitos de nuestra vida cotidiana.

En años recientes, los avances tecnológicos han provocado el surgimiento de nuevas áreas de estudio enfocadas en la gestión de la información. La Inteligencia Artificial (IA) es un área que representa un campo interesante y en constante crecimiento derivado del gran volumen de datos generados diariamente y a su facilidad para proponer soluciones a problemáticas que hasta hace unos años parecían difíciles o imposibles de resolver.

Los esfuerzos por parte de los gobiernos y autoridades de cada uno de los grupos y comunidades que conforman a la sociedad para garantizar la seguridad de sus individuos ha dado lugar a la adopción de nuevas tecnologías con la finalidad de contar con protocolos de seguridad más robustos. Es por esta razón que cada vez es más común encontrar la incorporación de tecnologías de última generación en sistemas para manejar el control de acceso a las instalaciones, verificar la identidad de los integrantes de su comunidad o simplemente garantizar un monitoreo continuo que permita la detección oportuna de anomalías y prevención de delitos.

La visión por computadora es un subcampo de la inteligencia artificial que ha tomado relevancia en los últimos años gracias a sus casos de uso que han permitido automatizar en gran de medida la forma en la que se proponen y desarrollan soluciones en materia de seguridad y vigilancia. De aquí se desprende el reconocimiento facial, una rama fundamental de la visión por computadora que permite identificar o verificar la identidad de un individuo a partir de sus características faciales en imágenes o vídeos.
Este campo ha visto cambios significativos impulsados en gran medida por el desarrollo del aprendizaje profundo, especialmente las redes neuronales convolucionales CNN.

Teniendo todo esto en mente es que se presenta el siguiente estado del arte, el cual proporciona una revisión breve sobre algunas técnicas usadas para llevar a cabo el reconocimiento facial y la detección de rostros, un paso importante dentro del flujo de trabajo de cualquier sistema de reconocimiento facial. Además, se prestara especial atención a los proyectos que integran estas tecnologías en un contexto estudiantil/académico, así como los desafíos y limitaciones de este campo de estudio.
